\documentclass[a4paper,11pt]{article}
\usepackage[utf8]{inputenc}
\usepackage[T1]{fontenc}
\usepackage[brazil]{babel}
\usepackage{amsmath}
\usepackage{graphicx}
\usepackage{booktabs}
\usepackage[margin=2.5cm]{geometry}
\usepackage{hyperref}

\title{MAC0425 - EP1: Agente de Busca para o Pac-Man}
\author{Leonardo Heidi Almeida Murakami - 11260186}
\date{\today}

\begin{document}

\maketitle

\section{Estatísticas das Buscas}

\subsection{Cenário 1 - Encontrando a comida num unico ponto}

\begin{table}[h]
\centering
\caption{Desempenho dos Algoritmos}
\begin{tabular}{llrr}
\toprule
\textbf{Algoritmo} & \textbf{Labirinto} & \textbf{Custo} & \textbf{Nós Expandidos} \\
\midrule
DFS & tinyMaze & 10 & 14 \\
    & mediumMaze & 130 & 144 \\
    & bigMaze & 210 & 390 \\
\midrule
BFS & tinyMaze & 8 & 15 \\
    & mediumMaze & 68 & 269 \\
    & bigMaze & 210 & 620 \\
\midrule
UCS & mediumMaze & 68 & 269 \\
    & mediumDottedMaze (StayEast) & 1 & 186 \\
    & mediumScaryMaze (StayWest) & 68719479864 & 108 \\
\midrule
A* (Manhattan) & tinyMaze & 8 & 14 \\
               & mediumMaze & 68 & 221 \\
               & bigMaze & 210 & 549 \\
\midrule
ID-DFS & tinyMaze & 8 & 58 \\
       & mediumMaze & 68 & 15429 \\
       & bigMaze & 210 & 60312 \\
\bottomrule
\end{tabular}
\end{table}

\begin{table}[h]
\centering
\caption{Desempenho do LRTA* no mediumMaze (Heurísticas Nula e Manhattan)}
\begin{tabular}{lrrr|rrr}
\toprule
 & \multicolumn{3}{c}{\textbf{Heurística Nula}} & \multicolumn{3}{c}{\textbf{Heurística Manhattan}} \\
\cmidrule(lr){2-4} \cmidrule(lr){5-7}
\textbf{Trials} & \textbf{Custo} & \textbf{Expandidos} & \textbf{H($s_0$)} & \textbf{Custo} & \textbf{Expandidos} & \textbf{H($s_0$)} \\
\midrule
10  & 76  & 3658  & 16 & 74 & 2130  & 56 \\
30  & 130 & 7486  & 30 & 68 & 5168  & 66 \\
50  & 110 & 9950  & 39 & 68 & 6580  & 68 \\
100 & 74  & 15034 & 59 & 68 & 9980  & 68 \\
150 & 68  & 18874 & 68 & 68 & 13380 & 68 \\
\bottomrule
\end{tabular}
\end{table}

\subsection{Cenário 2: Encontrando comida nos Cantos do Labirinto}

\begin{table}[h!]
\centering
\caption{Desempenho dos Algoritmos no Cenário 2 (Cantos)}
\begin{tabular}{llrr}
\toprule
\textbf{Algoritmo} & \textbf{Labirinto} & \textbf{Custo} & \textbf{Nós Expandidos} \\
\midrule
BFS & tinyCorners & 28 & 252 \\
    & mediumCorners & 106 & 1966 \\
    & bigCorners & 162 & 7949 \\
\midrule
A* (Corners) & tinyCorners & 28 & 154 \\
             & mediumCorners & 106 & 692 \\
             & bigCorners & 162 & 1725 \\
\bottomrule
\end{tabular}
\end{table}

\section{Discussão Comparativa}

O algoritmo \textbf{DFS} apresentou uma execução rápida e com baixo uso de memória nos ambientes maiores. No entanto, a principal desvantagem observada foi a incapacidade de garantir um caminho ótimo: no \textit{mediumMaze}, encontrou uma solução de custo 130, enquanto o custo ótimo é 68. 

Em contrapartida, a \textbf{BFS} explora o grafo de forma exaustiva (em largura), o que garante um plano de custo mínimo (ótimo de 68 no \textit{mediumMaze}). Essa exploração exaustiva aumenta consideravelmente a quantidade de nós expandidos em contraste com a DFS e, por consequência, demanda muito mais memória.

A \textbf{UCS} lida adequadamente com cenários em que as etapas têm custos desiguais. No \textit{mediumDottedMaze} (StayEast), o agente encontrou um caminho ótimo com custo 1 ao priorizar andar para a direita, em que o custo diminuía. No \textit{mediumScaryMaze} (StayWest), em que ir para a direita aumentava o custo exponencialmente ($2^x$), a UCS atenuou esse grande gasto evitando essas localizações, avaliando a função de custo dinamicamente em vez de apenas passos simples.

O algoritmo \textbf{A*} parece a escolha ideal de equilíbrio entre exploração e caminho ótimo: em ambos cenários, ele garantiu o caminho ótimo alcançando os mesmos custos retornados pelas cegas minimiais (BFS, UCS), porém reduzindo muito o número de expansões do BFS quando foi guiado por boas heurísticas, poupando memória.

O algoritmo \textbf{ID-DFS} alcançou o mesmo custo da DFS com o tamanho ótimo da BFS. Para conseguir isso sem utilizar uma memória extensa de estados enfileirados, foi preciso gerar as camadas da árvore repetidas vezes a cada profundidade investigada. Esse trade-off em tempo e limite resultou no maior número geral de nós expandidos em ambientes muito extensos, como o \textit{bigMaze} (mais de 60 mil).

Por fim, o \textbf{LRTA*} ilustrou, ao longo das tentativas repetidas (\textit{trials}) no \textit{mediumMaze}, como o agente calibra a estimativa. Usando a heurística nula, ele gradativamente aproximou o custo da intuição $H(s_0)$ próximo da verdade (~68) pelo processo on-line, em várias iterações com custo oscilando antes de ser constante. Com a \textbf{Manhattan Heuristic}, essa curva estabiliza quase que prontamente nas explorações iniciais por começar as buscas muito perto da melhor direção verdadeira ao final para os objetivos.

\section{Resposta das Questões}

\textbf{Questão 1:} Na DFS, a ordem de avaliação dos sucessores (Norte, Sul, Leste, Oeste) define o comportamento padrão de exploração do agente. Se o algoritmo não registrasse os estados já visitados (como ocorre em uma busca em árvore pura), o agente ficaria preso em ciclos infinitos. Em vez de avançar, ele andaria de um lado para o outro entre posições adjacentes, repetindo o mesmo caminho de ida e volta indefinidamente até estourar o limite de profundidade ou a memória do sistema.

\textbf{Questão 2:} A busca BFS sempre garante um caminho com o menor número de passos (custo mínimo, já que $1 \text{ passo} = \text{custo } 1$). Se o simulador alterar a ordem de avaliação dos sucessores, a BFS ainda encontrará uma rota otimizada. A única diferença é que ela poderá retornar um caminho alternativo, mas que possui exatamente o mesmo tamanho e custo total, garantindo a solução ótima.

\textbf{Questão 3:} No cenário \textit{StayEast}, ir para a esquerda tem um custo alto, enquanto caminhar para a direita diminui o custo. Por isso, a UCS (Busca de Custo Uniforme) prefere rotas que explorem mais o lado direito, minimizando o custo total acumulado. Por outro lado, no cenário \textit{StayWest}, andar para a direita gera uma penalidade exponencial ($2^x$). Nesse caso, a UCS muda a estratégia e passa a contornar o mapa pela esquerda, buscando evitar a todo custo os passos que geram multas tão altas.

\textbf{Questão 4:} A principal vantagem do A* em relação à UCS é que ele não considera apenas o custo real do caminho já percorrido ($g(s)$), mas também usa uma estimativa (heurística) de quão perto o agente está do objetivo ($h(s)$). Para que o A* sempre encontre o melhor caminho, essa heurística precisa ser consistente. Isso quer dizer que a distância estimada de um ponto até o de destino nunca deve criar atalhos irreais maiores que o custo real de locomoção.

\textbf{Questão 5:} No algoritmo LRTA*, o agente atualiza seu conhecimento sobre o labirinto de forma contínua enquanto se movimenta, ajustando suas estimativas com base na experiência (aprendizado on-line). Quando usamos a heurística de Manhattan, o agente já entra em campo com uma boa noção espacial. Isso significa que ele comete bem menos erros e estabiliza no caminho ótimo muito mais rápido. É um comportamento bem diferente do que acontece com a heurística nula, na qual o agente dá "pulos às cegas" nas primeiras tentativas até começar a memorizar o melhor trajeto.

\section{Melhorias e Dificuldades}

A principal dificuldade no desenvolvimento foi no problema dos cantos, em que foi preciso encontrar uma representação de estado que fosse leve e eficiente. Em vez de passar instâncias completas e pesadas da classe \textit{GameState} para cada nó da busca, optei por utilizar tuplas contendo apenas as informações essenciais (como a posição do agente e os cantos restantes). Essa escolha evita um alto consumo de memória e garante um ganho substancial de desempenho.

\end{document}
